% Options for packages loaded elsewhere
\PassOptionsToPackage{unicode}{hyperref}
\PassOptionsToPackage{hyphens}{url}
%
\documentclass[
]{book}
\usepackage{amsmath,amssymb}
\usepackage{iftex}
\ifPDFTeX
  \usepackage[T1]{fontenc}
  \usepackage[utf8]{inputenc}
  \usepackage{textcomp} % provide euro and other symbols
\else % if luatex or xetex
  \usepackage{unicode-math} % this also loads fontspec
  \defaultfontfeatures{Scale=MatchLowercase}
  \defaultfontfeatures[\rmfamily]{Ligatures=TeX,Scale=1}
\fi
\usepackage{lmodern}
\ifPDFTeX\else
  % xetex/luatex font selection
\fi
% Use upquote if available, for straight quotes in verbatim environments
\IfFileExists{upquote.sty}{\usepackage{upquote}}{}
\IfFileExists{microtype.sty}{% use microtype if available
  \usepackage[]{microtype}
  \UseMicrotypeSet[protrusion]{basicmath} % disable protrusion for tt fonts
}{}
\makeatletter
\@ifundefined{KOMAClassName}{% if non-KOMA class
  \IfFileExists{parskip.sty}{%
    \usepackage{parskip}
  }{% else
    \setlength{\parindent}{0pt}
    \setlength{\parskip}{6pt plus 2pt minus 1pt}}
}{% if KOMA class
  \KOMAoptions{parskip=half}}
\makeatother
\usepackage{xcolor}
\usepackage{color}
\usepackage{fancyvrb}
\newcommand{\VerbBar}{|}
\newcommand{\VERB}{\Verb[commandchars=\\\{\}]}
\DefineVerbatimEnvironment{Highlighting}{Verbatim}{commandchars=\\\{\}}
% Add ',fontsize=\small' for more characters per line
\usepackage{framed}
\definecolor{shadecolor}{RGB}{248,248,248}
\newenvironment{Shaded}{\begin{snugshade}}{\end{snugshade}}
\newcommand{\AlertTok}[1]{\textcolor[rgb]{0.94,0.16,0.16}{#1}}
\newcommand{\AnnotationTok}[1]{\textcolor[rgb]{0.56,0.35,0.01}{\textbf{\textit{#1}}}}
\newcommand{\AttributeTok}[1]{\textcolor[rgb]{0.13,0.29,0.53}{#1}}
\newcommand{\BaseNTok}[1]{\textcolor[rgb]{0.00,0.00,0.81}{#1}}
\newcommand{\BuiltInTok}[1]{#1}
\newcommand{\CharTok}[1]{\textcolor[rgb]{0.31,0.60,0.02}{#1}}
\newcommand{\CommentTok}[1]{\textcolor[rgb]{0.56,0.35,0.01}{\textit{#1}}}
\newcommand{\CommentVarTok}[1]{\textcolor[rgb]{0.56,0.35,0.01}{\textbf{\textit{#1}}}}
\newcommand{\ConstantTok}[1]{\textcolor[rgb]{0.56,0.35,0.01}{#1}}
\newcommand{\ControlFlowTok}[1]{\textcolor[rgb]{0.13,0.29,0.53}{\textbf{#1}}}
\newcommand{\DataTypeTok}[1]{\textcolor[rgb]{0.13,0.29,0.53}{#1}}
\newcommand{\DecValTok}[1]{\textcolor[rgb]{0.00,0.00,0.81}{#1}}
\newcommand{\DocumentationTok}[1]{\textcolor[rgb]{0.56,0.35,0.01}{\textbf{\textit{#1}}}}
\newcommand{\ErrorTok}[1]{\textcolor[rgb]{0.64,0.00,0.00}{\textbf{#1}}}
\newcommand{\ExtensionTok}[1]{#1}
\newcommand{\FloatTok}[1]{\textcolor[rgb]{0.00,0.00,0.81}{#1}}
\newcommand{\FunctionTok}[1]{\textcolor[rgb]{0.13,0.29,0.53}{\textbf{#1}}}
\newcommand{\ImportTok}[1]{#1}
\newcommand{\InformationTok}[1]{\textcolor[rgb]{0.56,0.35,0.01}{\textbf{\textit{#1}}}}
\newcommand{\KeywordTok}[1]{\textcolor[rgb]{0.13,0.29,0.53}{\textbf{#1}}}
\newcommand{\NormalTok}[1]{#1}
\newcommand{\OperatorTok}[1]{\textcolor[rgb]{0.81,0.36,0.00}{\textbf{#1}}}
\newcommand{\OtherTok}[1]{\textcolor[rgb]{0.56,0.35,0.01}{#1}}
\newcommand{\PreprocessorTok}[1]{\textcolor[rgb]{0.56,0.35,0.01}{\textit{#1}}}
\newcommand{\RegionMarkerTok}[1]{#1}
\newcommand{\SpecialCharTok}[1]{\textcolor[rgb]{0.81,0.36,0.00}{\textbf{#1}}}
\newcommand{\SpecialStringTok}[1]{\textcolor[rgb]{0.31,0.60,0.02}{#1}}
\newcommand{\StringTok}[1]{\textcolor[rgb]{0.31,0.60,0.02}{#1}}
\newcommand{\VariableTok}[1]{\textcolor[rgb]{0.00,0.00,0.00}{#1}}
\newcommand{\VerbatimStringTok}[1]{\textcolor[rgb]{0.31,0.60,0.02}{#1}}
\newcommand{\WarningTok}[1]{\textcolor[rgb]{0.56,0.35,0.01}{\textbf{\textit{#1}}}}
\usepackage{longtable,booktabs,array}
\usepackage{calc} % for calculating minipage widths
% Correct order of tables after \paragraph or \subparagraph
\usepackage{etoolbox}
\makeatletter
\patchcmd\longtable{\par}{\if@noskipsec\mbox{}\fi\par}{}{}
\makeatother
% Allow footnotes in longtable head/foot
\IfFileExists{footnotehyper.sty}{\usepackage{footnotehyper}}{\usepackage{footnote}}
\makesavenoteenv{longtable}
\usepackage{graphicx}
\makeatletter
\def\maxwidth{\ifdim\Gin@nat@width>\linewidth\linewidth\else\Gin@nat@width\fi}
\def\maxheight{\ifdim\Gin@nat@height>\textheight\textheight\else\Gin@nat@height\fi}
\makeatother
% Scale images if necessary, so that they will not overflow the page
% margins by default, and it is still possible to overwrite the defaults
% using explicit options in \includegraphics[width, height, ...]{}
\setkeys{Gin}{width=\maxwidth,height=\maxheight,keepaspectratio}
% Set default figure placement to htbp
\makeatletter
\def\fps@figure{htbp}
\makeatother
\setlength{\emergencystretch}{3em} % prevent overfull lines
\providecommand{\tightlist}{%
  \setlength{\itemsep}{0pt}\setlength{\parskip}{0pt}}
\setcounter{secnumdepth}{5}
\ifLuaTeX
\usepackage[bidi=basic]{babel}
\else
\usepackage[bidi=default]{babel}
\fi
\babelprovide[main,import]{french}
% get rid of language-specific shorthands (see #6817):
\let\LanguageShortHands\languageshorthands
\def\languageshorthands#1{}
\usepackage{booktabs}
\usepackage{datetime}
\ifLuaTeX
  \usepackage{selnolig}  % disable illegal ligatures
\fi
\usepackage{bookmark}
\IfFileExists{xurl.sty}{\usepackage{xurl}}{} % add URL line breaks if available
\urlstyle{same}
\hypersetup{
  pdftitle={Analyse de l'impact de la pression anthropique sur la biodiversité dans le bassin du Congo},
  pdfauthor={Cica TOSSOU, Thibaut TOSSOU, Estève Alvin KASSA},
  pdflang={fr},
  hidelinks,
  pdfcreator={LaTeX via pandoc}}

\title{Analyse de l'impact de la pression anthropique sur la biodiversité dans le bassin du Congo}
\author{Cica TOSSOU, Thibaut TOSSOU, Estève Alvin KASSA}
\date{2025-04-17}

\begin{document}
\maketitle

{
\setcounter{tocdepth}{1}
\tableofcontents
}
\newpage

Ce document présente les résultats d'une étude sur l'a'impact de l'activité humaine sur une population de céphalophes dans cinq villages du bassin du Congo. Elle est réalisée dans le cadre de l'évaluation formative du temps 1, de la certification en analyse de données massives du Groupe des Ecoles Centrales en partenariaat avec Sèmè City au Bénin.

\newpage

\chapter{Résumé de l'étude}\label{ruxe9sumuxe9-de-luxe9tude}

L'étude révèle que l'abondance des céphalophes varie selon l'altitude, la couverture forestière et la distance aux activités humaines. C. dorsalis préfère les zones denses et isolées, C. silvicultor les forêts en altitude, tandis que P. maxwellii tolère mieux la proximité humaine. Ces résultats soulignent l'importance de stratégies différenciées de conservation selon les espèces.

\chapter{L'introduction}\label{lintroduction}

Les céphalophes sont des petits bovidés forestiers discrets, mais essentiels au fonctionnement des écosystèmes tropicaux. Leur présence et leur abondance sont influencées par de nombreux facteurs environnementaux et anthropiques, souvent complexes à démêler. Dans un contexte de fragmentation croissante des habitats et de pressions humaines accrues, comprendre les préférences écologiques de ces espèces devient crucial pour orienter efficacement les efforts de conservation.

Cette étude s'appuie sur des relevés de pièges photographiques menés dans une aire protégée de l'Afrique de l'Ouest. Elle vise à évaluer l'influence de variables environnementales telles que l'altitude, la couverture forestière et la proximité des activités humaines sur l'abondance de trois espèces de céphalophes : Cephalophus dorsalis, Cephalophus silvicultor et Philantomba maxwellii. En identifiant les préférences spécifiques de chaque espèce, l'étude apporte des éléments concrets pour mieux planifier la gestion des habitats forestiers.

\chapter{Méthodologie de l'étude}\label{muxe9thodologie-de-luxe9tude}

(Ajouter ici une introduction de la section méthodologie)

\chapter{Données brutes mises à disposition}\label{donnuxe9es-brutes-mises-uxe0-disposition}

Pour cette étude, nous avons reçu un fichier excel nommé \textbf{Data\_Biodiversity.xlsx}. Ce dernier agrège en 4 feuilles, les données quantitatives et qualitatives mises à notre disposition dans le cadre d cet étude :

\begin{itemize}
\tightlist
\item
  data\_espèces : nous y retrouvons, les détails concernant les différentes apparitions de céphalophes enrégistrées par les stations de collecte déployées sur les cinq villages ;
\item
  data\_stations : répertorie l'ensemble des installations de station sur le territoire. Cette feuille donne aussi, la nomenclature des noms des stations d'échantillonage, ainsi que des données importantes sur leur lieu d'implantation (relief, caractéristiques naturelles, distance aux routes, etc\ldots)
\item
  table\_information : dictionnaire de données de la base \textbf{Data\_Biodiversity.xlsx}
\item
  autres\_infos : s'y trouvent des informations importantes pour identifier les différentes espèces et la superficie des villages échantillonés.
\end{itemize}

\textbf{Data\_Biodiversity.xlsx} est propre, bien structuré, et ne contient pas de données manquantes au début de l'étude.

\chapter{Préparation des données}\label{pruxe9paration-des-donnuxe9es}

Nous avons extrait les informations importantes de la base de données, en utilisant le package openxlsx. C'est celui qui offrait la meilleure flexibilité dans le choix des cellules à extraire de la base fournie.

\begin{Shaded}
\begin{Highlighting}[]
\CommentTok{\#Activer la librairie pour l\textquotesingle{}utiliser}
\FunctionTok{library}\NormalTok{(openxlsx)}

\CommentTok{\#Extraction des données importantes}
\NormalTok{data\_especes }\OtherTok{\textless{}{-}} \FunctionTok{read.xlsx}\NormalTok{(}\StringTok{"C:/Users/HP/OneDrive {-} gouvbj/MVDC/Cours/Projects/Biodiversity/data/raw/Data\_Biodiversity.xlsx"}\NormalTok{,}
                             \AttributeTok{sheet =} \DecValTok{1}\NormalTok{,}
                             \AttributeTok{startRow =} \DecValTok{1}\NormalTok{,}
                             \AttributeTok{colNames =} \ConstantTok{TRUE}\NormalTok{,}
                             \AttributeTok{rowNames =} \ConstantTok{FALSE}\NormalTok{,}
                             \AttributeTok{detectDates =} \ConstantTok{TRUE}\NormalTok{,}
                             \AttributeTok{sep.names =} \StringTok{"\_"}\NormalTok{,}
                             \AttributeTok{na.strings =} \StringTok{"NA"}\NormalTok{)}

\FunctionTok{str}\NormalTok{(data\_especes)}
\end{Highlighting}
\end{Shaded}

\begin{verbatim}
## 'data.frame':    2545 obs. of  5 variables:
##  $ Villages                  : chr  "Sekom" "Sekom" "Sekom" "Sekom" ...
##  $ Stations_d'échantillonnage: chr  "Sekom01A" "Sekom01A" "Sekom01A" "Sekom01A" ...
##  $ Espèces                   : chr  "P. monticola" "P. monticola" "C. callipygus" "P. monticola" ...
##  $ Heure_d'observation       : chr  "0.52638888888888891" "15:13" "17:32" "13:13" ...
##  $ Dates_d'échantillonnage   : Date, format: "2014-06-04" "2014-06-10" ...
\end{verbatim}

\begin{Shaded}
\begin{Highlighting}[]
\CommentTok{\#Extraction des données importantes}
\NormalTok{data\_especes }\OtherTok{\textless{}{-}} \FunctionTok{read.xlsx}\NormalTok{(}\StringTok{"C:/Users/HP/OneDrive {-} gouvbj/MVDC/Cours/Projects/Biodiversity/data/raw/Data\_Biodiversity.xlsx"}\NormalTok{,}
                             \AttributeTok{sheet =} \DecValTok{1}\NormalTok{,}
                             \AttributeTok{startRow =} \DecValTok{1}\NormalTok{,}
                             \AttributeTok{colNames =} \ConstantTok{TRUE}\NormalTok{,}
                             \AttributeTok{rowNames =} \ConstantTok{FALSE}\NormalTok{,}
                             \AttributeTok{detectDates =} \ConstantTok{TRUE}\NormalTok{,}
                             \AttributeTok{sep.names =} \StringTok{"\_"}\NormalTok{,}
                             \AttributeTok{na.strings =} \StringTok{"NA"}\NormalTok{)}

\FunctionTok{str}\NormalTok{(data\_especes)}
\end{Highlighting}
\end{Shaded}

\begin{verbatim}
## 'data.frame':    2545 obs. of  5 variables:
##  $ Villages                  : chr  "Sekom" "Sekom" "Sekom" "Sekom" ...
##  $ Stations_d'échantillonnage: chr  "Sekom01A" "Sekom01A" "Sekom01A" "Sekom01A" ...
##  $ Espèces                   : chr  "P. monticola" "P. monticola" "C. callipygus" "P. monticola" ...
##  $ Heure_d'observation       : chr  "0.52638888888888891" "15:13" "17:32" "13:13" ...
##  $ Dates_d'échantillonnage   : Date, format: "2014-06-04" "2014-06-10" ...
\end{verbatim}

\chapter{Résultats de l'analyse exploratoire et statistique}\label{ruxe9sultats-de-lanalyse-exploratoire-et-statistique}

\section{Footnotes}\label{footnotes}

Footnotes are put inside the square brackets after a caret \texttt{\^{}{[}{]}}. Like this one \footnote{This is a footnote.}.

\section{Citations}\label{citations}

Reference items in your bibliography file(s) using \texttt{@key}.

\chapter{Discussions}\label{discussions}

\chapter{Conclusion}\label{conclusion}

\chapter{Parts}\label{parts}

\end{document}
